% GENERATED FILE. Edit canonical lessons, then run npm run generate:books.
% canonical-source-hash: fnv1a64:66bf6c7090e8a806
% canonical-lessons: ES-C568-historico, ES-C568-revolucion, ES-C568-monarca, ES-C568-reina, ES-C568-conquistar

\chapter{Reyes y revoluciones}
\label{ch:es-reyes-y-revoluciones}
\hlchaptermodality{568}

\begin{chapteropening}
\textbf{By the end of this chapter:} \emph{I can talk about history: monarchs and queens, conquests, revolutions and historic days.}

Complete the last lesson of chapter 568: I can talk about history: monarchs and queens, conquests, revolutions and historic days.
\end{chapteropening}

\section[histórico]{histórico — historic; historical}

\label{lesson:ES-C568-historico}



\begin{quote}
Before the new one: say the Spanish for the soldier, then the Spanish for the army.
\end{quote}

\subsection*{You'll want to know: histórico}
\textbf{histórico} — \textquotedblleft{}historic; historical\textquotedblright{}.

\begin{quote}
Fue un día histórico para el país.
\end{quote}

\begin{quote}
\emph{It was a historic day for the country.}
\end{quote}

\emph{El centro histórico} — the old town.

\begin{cousinweb}
\emph{histórico} is a learned borrowing from Latin \emph{historicus}, which comes from Greek \emph{historikós} 'of history' (EtymDB), built on \emph{historía} 'an inquiry'.
\end{cousinweb}

\begin{grammarlens}[title={What you've built}]
One more word for telling, explaining and describing at B1.
\end{grammarlens}

\subsection*{Your turn}
\begin{itemize}
\raggedright
  \item \emph{Say it:} \emph{histórico}
  \item \emph{Say it:} \emph{histórico}, once more
  \item \emph{Say it:} a sentence of your own with \emph{histórico}
  \item \emph{Recall:} say the Spanish for the soldier, then the Spanish for the army, then say \emph{histórico} again
\end{itemize}

\begin{tcolorbox}[breakable,colback=teal!4,colframe=teal!35!black,title={Before you move on}]
What does histórico mean? (\textquotedblleft{}Historic; historical\textquotedblright{}.) Say it once more.
\end{tcolorbox}
\section[la revolución]{la revolución — the revolution}

\label{lesson:ES-C568-revolucion}



\begin{quote}
Before the new one: say the Spanish for the army, then the Spanish for historic.
\end{quote}

\subsection*{You'll want to know: la revolución}
\textbf{la revolución} — \textquotedblleft{}the revolution\textquotedblright{}.

\begin{quote}
La revolución cambió el país.
\end{quote}

\begin{quote}
\emph{The revolution changed the country.}
\end{quote}

\begin{cousinweb}
\emph{revolución} is a learned borrowing from Late Latin \emph{revolūtiō} 'a rolling back, a revolving' (EtymDB), built on \emph{re-} and \emph{volvere} 'to roll, to turn'.
\end{cousinweb}

\begin{grammarlens}[title={What you've built}]
One more word for telling, explaining and describing at B1.
\end{grammarlens}

\subsection*{Your turn}
\begin{itemize}
\raggedright
  \item \emph{Say it:} \emph{la revolución}
  \item \emph{Say it:} \emph{la revolución}, once more
  \item \emph{Say it:} a sentence of your own with \emph{la revolución}
  \item \emph{Recall:} say the Spanish for the army, then the Spanish for historic, then say \emph{la revolución} again
\end{itemize}

\begin{tcolorbox}[breakable,colback=teal!4,colframe=teal!35!black,title={Before you move on}]
What does la revolución mean? (\textquotedblleft{}The revolution\textquotedblright{}.) Say it once more.
\end{tcolorbox}
\section[el monarca]{el monarca — the monarch}

\label{lesson:ES-C568-monarca}



\begin{quote}
Before the new one: say the Spanish for historic, then the Spanish for the revolution.
\end{quote}

\subsection*{You'll want to know: el monarca}
\textbf{el monarca} — \textquotedblleft{}the monarch\textquotedblright{}.

\begin{quote}
El monarca visitó la ciudad en 1990.
\end{quote}

\begin{quote}
\emph{The monarch visited the city in 1990.}
\end{quote}

Masculine although it ends in \emph{-a}. \emph{La monarquía} is the monarchy.

\begin{cousinweb}
\emph{monarca} is a learned borrowing from Late Latin \emph{monarcha}, which comes from Greek \emph{monárchēs} 'one who rules alone' (EtymDB), built on \emph{mónos} 'alone' and \emph{árchein} 'to rule'.
\end{cousinweb}

\begin{grammarlens}[title={What you've built}]
One more word for telling, explaining and describing at B1.
\end{grammarlens}

\subsection*{Your turn}
\begin{itemize}
\raggedright
  \item \emph{Say it:} \emph{el monarca}
  \item \emph{Say it:} \emph{el monarca}, once more
  \item \emph{Say it:} a sentence of your own with \emph{el monarca}
  \item \emph{Recall:} say the Spanish for historic, then the Spanish for the revolution, then say \emph{el monarca} again
\end{itemize}

\begin{tcolorbox}[breakable,colback=teal!4,colframe=teal!35!black,title={Before you move on}]
What does el monarca mean? (\textquotedblleft{}The monarch\textquotedblright{}.) Say it once more.
\end{tcolorbox}
\section[la reina]{la reina — the queen}

\label{lesson:ES-C568-reina}



\begin{quote}
Before the new one: say the Spanish for the revolution, then the Spanish for the monarch.
\end{quote}

\subsection*{You'll want to know: la reina}
\textbf{la reina} — \textquotedblleft{}the queen\textquotedblright{}.

\begin{quote}
La reina inauguró el hospital.
\end{quote}

\begin{quote}
\emph{The queen opened the hospital.}
\end{quote}

\begin{cousinweb}
\emph{reina} is inherited from Latin \emph{rēgīna} 'queen' (EtymDB), through old Spanish \emph{reyna}. \emph{Rēgīna} is built on \emph{rēx} 'king'.
\end{cousinweb}

\begin{grammarlens}[title={What you've built}]
One more word for telling, explaining and describing at B1.
\end{grammarlens}

\subsection*{Your turn}
\begin{itemize}
\raggedright
  \item \emph{Say it:} \emph{la reina}
  \item \emph{Say it:} \emph{la reina}, once more
  \item \emph{Say it:} a sentence of your own with \emph{la reina}
  \item \emph{Recall:} say the Spanish for the revolution, then the Spanish for the monarch, then say \emph{la reina} again
\end{itemize}

\begin{tcolorbox}[breakable,colback=teal!4,colframe=teal!35!black,title={Before you move on}]
What does la reina mean? (\textquotedblleft{}The queen\textquotedblright{}.) Say it once more.
\end{tcolorbox}
\section[conquistar]{conquistar — to conquer; to win over}

\label{lesson:ES-C568-conquistar}



\begin{quote}
Before the new one: say the Spanish for the monarch, then the Spanish for the queen.
\end{quote}

\subsection*{You'll want to know: conquistar}
\textbf{conquistar} — \textquotedblleft{}to conquer; to win over\textquotedblright{}.

\begin{quote}
Los romanos conquistaron la península.
\end{quote}

\begin{quote}
\emph{The Romans conquered the peninsula.}
\end{quote}

It also means to win someone's heart: \emph{La conquistó con su humor.}

\begin{cousinweb}
\emph{conquistar} is usually traced to spoken Latin \emph{*conquīstāre}, built on \emph{conquīsītus}, the participle of \emph{conquīrere} 'to seek out, to bring together' (Lewis \& Short), from \emph{quaerere} 'to seek'.
\end{cousinweb}

\begin{grammarlens}[title={What you've built}]
One more word for telling, explaining and describing at B1.
\end{grammarlens}

\subsection*{Your turn}
\begin{itemize}
\raggedright
  \item \emph{Say it:} \emph{conquistar}
  \item \emph{Say it:} \emph{conquistar}, once more
  \item \emph{Say it:} a sentence of your own with \emph{conquistar}
  \item \emph{Recall:} say the Spanish for the monarch, then the Spanish for the queen, then say \emph{conquistar} again
\end{itemize}

\begin{tcolorbox}[breakable,colback=teal!4,colframe=teal!35!black,title={Before you move on}]
What does conquistar mean? (\textquotedblleft{}To conquer; to win over\textquotedblright{}.) Say it once more.
\end{tcolorbox}
